\documentclass[12pt,a4paper]{article}

% ============================================================
%  CODIFICAÇÃO E IDIOMA
% ============================================================
\usepackage[utf8]{inputenc}
\usepackage[T1]{fontenc}
\usepackage[brazil]{babel}
% Trocar "brazil" por "english" se decidir submeter em inglês
% (Principia aceita PT/EN/FR/ES).

% ============================================================
%  MATEMÁTICA E AMBIENTES FORMAIS (estilo Suppes/Bourbaki)
% ============================================================
\usepackage{amsmath,amssymb,amsthm}

\theoremstyle{definition}
\newtheorem{definition}{Definição}[section]
\newtheorem{axiom}{Axioma}[section]
\theoremstyle{plain}
\newtheorem{proposition}{Proposição}[section]
\theoremstyle{remark}
\newtheorem*{example}{Exemplo}

% ============================================================
%  LAYOUT
%  Valores abaixo são convenção acadêmica genérica, não a norma
%  confirmada da chamada Principia/Transversal do I CBFC — checar
%  o edital específico antes de formatar a versão final.
% ============================================================
\usepackage[margin=2.5cm]{geometry}
\usepackage{setspace}
\doublespacing

% ============================================================
%  CITAÇÕES — autor-data (padrão dominante em filosofia)
% ============================================================
\usepackage[authoryear,round]{natbib}

\usepackage{hyperref}
\hypersetup{colorlinks=true, linkcolor=blue, citecolor=blue, urlcolor=blue}

% ============================================================
%  METADADOS
%  Principia exige revisão cega: remova nome/afiliação/agradecimentos
%  do corpo do PDF de avaliação; essas informações vão em cadastro
%  separado no sistema de submissão, não no manuscrito.
% ============================================================
\title{Honestidade Axiomática como um Critério Normativo e Pluralista de Segunda Ordem para Modelos em Teorias Científicas}
\author{[nome removido para revisão cega]}
\date{}

\begin{document}
	\maketitle
	
	\begin{abstract}
		% Resumo em português, 150–250 palavras. Escrever por último,
		% depois do rascunho pronto — não antes.
	\end{abstract}
	
	\selectlanguage{english}
	\begin{abstract}
		% English abstract. Praticamente obrigatório em periódico
		% internacional mesmo submetendo o corpo em português.
	\end{abstract}
	\selectlanguage{brazil}
	
	\noindent\textbf{Palavras-chave:} honestidade axiomática; adequação estatística; filosofia da ciência; séries temporais fuzzy.
	
	% ================================================================
	\section{Introdução}
	% PARÁGRAFO ZERO — escrever isto primeiro, antes de qualquer outra
	% seção. É a lacuna que o I CBFC apontou: onde honestidade axiomática
	% se situa face a realismo / anti-realismo / positivismo. Declare a
	% posição (normativa, não descritiva; estruturalista via
	% Suppes/Krause) explicitamente, não deixe implícita.
	
	
	
	
	% ================================================================
	\section{Honestidade Axiomática}
	% Definição formal + grade 2x2 (adequação x honestidade) +
	% predicado normativo como a contribuição central.
	
	\begin{definition}[Honestidade Axiomática]
		
	\end{definition}
	
	
	% ================================================================
	\section{Caso Canônico: Kepler--Ptolomeu}
	% \citep{spanos2007} carrega o peso teórico sem a carga técnica de
	% regressão que afastou a audiência no congresso.
	
	
	% ================================================================
	\section{Casos de Apoio}
	% Porto de Santos (PWFTS diff(12)) e Kazi/Belgrado \citep{kazi2023}
	% como exemplos de R^2 ao lado do caso canônico — não mais como
	% estudo de caso central isolado.
	
	
	% ================================================================
	\section{Formalização via Suppes}
	% Passada de rigor sobre a definição da Seção 2, usando o critério
	% de invariância/meaningfulness (Suppes, \emph{Models and Methods},
	% cap. 3). Escrever isto por ÚLTIMO, depois do rascunho completo —
	% é reaplicação do critério já extraído, não pré-leitura.
	
	
	% ================================================================
	\section{Considerações Finais}
	
	
	% ================================================================
	% BIBLIOGRAFIA
	% Chaves já usadas acima pressupõem um referencias.bib com, no
	% mínimo: spanos2007, kazi2023, mayospanos2004, grangernewbold1974,
	% krause2022, nielsen2020, hyndman2025. Ainda não existe — avise se
	% quiser o stub gerado.
	% ================================================================
	\bibliographystyle{apalike}
	\bibliography{referencias}
	
\end{document}